% References. Entries checked against DOI / arXiv on 27 Sept (docs/thesis/references.md) or on 2 Oct (this report).
\begin{thebibliography}{99}\small

\bibitem{alonzo2022beyond} O.~Alonzo, H.~V. Shin, D.~Li. Beyond subtitles: Captioning and visualizing non-speech sounds to improve accessibility of user-generated videos. \emph{ASSETS}, 2022. doi:10.1145/3517428.3544808.
\bibitem{alhamoud2025negbench} K.~Alhamoud et al. Vision-language models do not understand negation (NegBench). \emph{CVPR}, 2025. arXiv:2501.09425.
\bibitem{bai2025qwen25vl} Qwen Team (S.~Bai et al.). Qwen2.5-VL technical report. arXiv:2502.13923, 2025.
\bibitem{bfl2024flux} Black Forest Labs. FLUX.1 [schnell]. Model release, 2024. \url{https://huggingface.co/black-forest-labs/FLUX.1-schnell}.
\bibitem{bilen2020psds} \c{C}.~Bilen, G.~Ferroni, F.~Tuveri, J.~Azcarreta, S.~Krstulovi\'c. A framework for the robust evaluation of sound event detection. \emph{ICASSP}, 2020. arXiv:1910.08440.
\bibitem{biner2024sonicdiffusion} B.~C. Biner, F.~M. Sofian, U.~B. Karaka\c{s}, D.~Ceylan, E.~Erdem, A.~Erdem. SonicDiffusion: Audio-driven image generation and editing with pretrained diffusion models. arXiv:2405.00878, 2024.
\bibitem{chen2020vggsound} H.~Chen, W.~Xie, A.~Vedaldi, A.~Zisserman. VGGSound: A large-scale audio-visual dataset. \emph{ICASSP}, 2020. arXiv:2004.14368.
\bibitem{chen2021hardway} H.~Chen, W.~Xie, T.~Afouras, A.~Nagrani, A.~Vedaldi, A.~Zisserman. Localizing visual sounds the hard way. \emph{CVPR}, 2021. arXiv:2104.02691.
\bibitem{chen2023beats} S.~Chen, Y.~Wu, C.~Wang, S.~Liu, D.~Tompkins, Z.~Chen, F.~Wei. BEATs: Audio pre-training with acoustic tokenizers. \emph{ICML}, 2023. arXiv:2212.09058.
\bibitem{chu2023qwenaudio} Y.~Chu, J.~Xu, X.~Zhou, Q.~Yang, S.~Zhang, Z.~Yan, C.~Zhou, J.~Zhou. Qwen-Audio: Advancing universal audio understanding via unified large-scale audio-language models. arXiv:2311.07919, 2023.
\bibitem{chu2024qwen2audio} Y.~Chu et al. Qwen2-Audio technical report. arXiv:2407.10759, 2024.
\bibitem{dcmp} Described and Captioned Media Program. Captioning Key: Sound effects and music. \url{https://dcmp.org/learn/captioningkey}, retrieved 27 Sept 2026.
\bibitem{delacerdapataca2024royale} C.~de~Lacerda~Pataca, S.~Hassan, N.~Tinker, R.~L. Peiris, M.~Huenerfauth. Caption Royale: Exploring the design space of affective captions from the perspective of deaf and hard-of-hearing individuals. \emph{CHI}, 2024. doi:10.1145/3613904.3642258.
\bibitem{drummond2006costcurves} C.~Drummond, R.~C. Holte. Cost curves: An improved method for visualizing classifier performance. \emph{Machine Learning} 65(1):95--130, 2006.
\bibitem{elizalde2023clap} B.~Elizalde, S.~Deshmukh, M.~Al~Ismail, H.~Wang. CLAP: Learning audio concepts from natural language supervision. \emph{ICASSP}, 2023.
\bibitem{ferroni2021dcase} G.~Ferroni, N.~Turpault, J.~Azcarreta, F.~Tuveri, R.~Serizel, \c{C}.~Bilen, S.~Krstulovi\'c. Improving sound event detection metrics: Insights from DCASE 2020. \emph{ICASSP}, 2021. arXiv:2010.13648.
\bibitem{ge2026audiocanvas} D.~Ge, S.~Liu, C.~Gong, X.-L. Zhang, C.~Zhang, X.~Li. Towards expressive and faithful audio-to-image generation: A unified multimodal dataset and synthesis framework (AudioCanvas). \emph{ACM MM}, 2026. arXiv:2608.09529.
\bibitem{gemmeke2017audioset} J.~F. Gemmeke, D.~P.~W. Ellis, D.~Freedman, A.~Jansen, W.~Lawrence, R.~C. Moore, M.~Plakal, M.~Ritter. Audio Set: An ontology and human-labeled dataset for audio events. \emph{ICASSP}, 2017.
\bibitem{geng2023unav100} T.~Geng, T.~Wang, J.~Duan, R.~Cong, F.~Zheng. Dense-localizing audio-visual events in untrimmed videos: A large-scale benchmark and baseline. \emph{CVPR}, 2023. arXiv:2303.12930.
\bibitem{gong2021ast} Y.~Gong, Y.-A. Chung, J.~Glass. AST: Audio spectrogram transformer. \emph{Interspeech}, 2021. arXiv:2104.01778.
\bibitem{hai2025flexsed} J.~Hai, H.~Wang, W.~Guo, M.~Elhilali. FlexSED: Towards open-vocabulary sound event detection. \emph{WASPAA}, 2025. arXiv:2509.18606.
\bibitem{hershey2021strong} S.~Hershey, D.~P.~W. Ellis, E.~Fonseca, A.~Jansen, C.~Liu, R.~C. Moore, M.~Plakal. The benefit of temporally-strong labels in audio event classification. \emph{ICASSP}, 2021. arXiv:2105.07031.
\bibitem{homesound} D.~Jain, K.~Mack, A.~Amrous, M.~Wright, S.~Goodman, L.~Findlater, J.~E. Froehlich. HomeSound: An iterative field deployment of an in-home sound awareness system for deaf or hard of hearing users. \emph{CHI}, 2020. doi:10.1145/3313831.3376758.
\bibitem{jain2015hmd} D.~Jain, L.~Findlater, J.~Gilkeson, B.~Holland, R.~Duraiswami, D.~Zotkin, C.~Vogler, J.~E. Froehlich. Head-mounted display visualizations to support sound awareness for the deaf and hard of hearing. \emph{CHI}, 2015.
\bibitem{jain2019home} D.~Jain, A.~Lin, R.~Guttman, M.~Amalachandran, A.~Zeng, L.~Findlater, J.~Froehlich. Exploring sound awareness in the home for people who are deaf or hard of hearing. \emph{CHI}, 2019.
\bibitem{juanola2025critical} X.~Juanola, G.~Haro, M.~Fuentes. A critical assessment of visual sound source localization models including negative audio. \emph{ICASSP}, 2025. arXiv:2410.01020.
\bibitem{juanola2025silence} X.~Juanola, G.~Morais, M.~Fuentes, G.~Haro. Learning from silence and noise for visual sound source localization. \emph{BMVC}, 2025. arXiv:2508.21761.
\bibitem{kong2020panns} Q.~Kong, Y.~Cao, T.~Iqbal, Y.~Wang, W.~Wang, M.~D. Plumbley. PANNs: Large-scale pretrained audio neural networks for audio pattern recognition. \emph{IEEE/ACM TASLP} 28:2880--2894, 2020.
\bibitem{liu2023llava} H.~Liu, C.~Li, Q.~Wu, Y.~J. Lee. Visual instruction tuning. \emph{NeurIPS}, 2023. arXiv:2304.08485.
\bibitem{matthews2006nonspeech} T.~Matthews, J.~Fong, F.~W.-L. Ho-Ching, J.~Mankoff. Evaluating non-speech sound visualizations for the deaf. \emph{Behaviour \& IT} 25(4):333--351, 2006.
\bibitem{may2024unspoken} L.~May, K.~Ohshiro, K.~Dang, S.~Sridhar, J.~Pai, M.~Fuentes, S.~Lee, M.~Cartwright. Unspoken sound: Identifying trends in non-speech audio captioning on YouTube. \emph{CHI}, 2024. doi:10.1145/3613904.3642162.
\bibitem{may2025choices} L.~May, M.~Clemens, K.~Dang, K.~Ohshiro, S.~Sridhar, P.~Wee, M.~Fuentes, S.~Lee, M.~Cartwright. ``Choices? That's the dream'': Challenges and opportunities in non-speech information closed-captioning. \emph{Frontiers in Computer Science} 7, 2025.
\bibitem{mesaros2016metrics} A.~Mesaros, T.~Heittola, T.~Virtanen. Metrics for polyphonic sound event detection. \emph{Applied Sciences} 6(6):162, 2016.
\bibitem{owens2018multisensory} A.~Owens, A.~A. Efros. Audio-visual scene analysis with self-supervised multisensory features. \emph{ECCV}, 2018. arXiv:1804.03641.
\bibitem{petermann2025seeingsound} D.~Petermann, M.~M. Kalayeh. Seeing Sound: Assembling sounds from visuals for audio-to-image generation. arXiv:2501.05413, 2025.
\bibitem{podell2023sdxl} D.~Podell et al. SDXL: Improving latent diffusion models for high-resolution image synthesis. arXiv:2307.01952, 2023.
\bibitem{prosodycaptions} C.~de~Lacerda~Pataca, M.~Watkins, R.~Peiris, S.~Lee, M.~Huenerfauth. Visualization of speech prosody and emotion in captions: Accessibility for deaf and hard-of-hearing users. \emph{CHI}, 2023. doi:10.1145/3544548.3581511.
\bibitem{protosound} D.~Jain et al. ProtoSound: A personalized and scalable sound recognition system for deaf and hard-of-hearing users. \emph{CHI}, 2022. arXiv:2202.11134.
\bibitem{revuelta2020eeg} P.~Revuelta, T.~Ortiz, M.~J. Luc\'{\i}a, B.~Ruiz, J.~M. S\'anchez-Pena. Limitations of standard accessible captioning of sounds and music for deaf and hard of hearing people: An EEG study. \emph{Frontiers in Integrative Neuroscience} 14, 2020.
\bibitem{sakshi2024mmau} S.~Sakshi et al. MMAU: A massive multi-task audio understanding and reasoning benchmark. arXiv:2410.19168, 2024.
\bibitem{senocak2018localize} A.~Senocak, T.-H. Oh, J.~Kim, M.-H. Yang, I.~S. Kweon. Learning to localize sound source in visual scenes. \emph{CVPR}, 2018.
\bibitem{senocak2023alignment} A.~Senocak, H.~Ryu, J.~Kim, T.-H. Oh, H.~Pfister, J.~S. Chung. Sound source localization is all about cross-modal alignment. \emph{ICCV}, 2023. arXiv:2309.10724.
\bibitem{soundwatch} D.~Jain, H.~Ngo, P.~Patel, S.~Goodman, L.~Findlater, J.~Froehlich. SoundWatch: Exploring smartwatch-based deep learning approaches to support sound awareness for deaf and hard of hearing users. \emph{ASSETS}, 2020.
\bibitem{sungbin2023sound2scene} K.~Sung-Bin, A.~Senocak, H.~Ha, A.~Owens, T.-H. Oh. Sound to visual scene generation by audio-to-visual latent alignment. \emph{CVPR}, 2023. arXiv:2303.17490.
\bibitem{tang2024salmonn} C.~Tang et al. SALMONN: Towards generic hearing abilities for large language models. \emph{ICLR}, 2024. arXiv:2310.13289.
\bibitem{tzinis2021audioscope} E.~Tzinis, S.~Wisdom, A.~Jansen, S.~Hershey, T.~Remez, D.~P.~W. Ellis, J.~R. Hershey. Into the wild with AudioScope: Unsupervised audio-visual separation of on-screen sounds. \emph{ICLR}, 2021. arXiv:2011.01143.
\bibitem{um2025objaware} S.~J. Um, D.~Kim, S.~Lee, J.~U. Kim. Object-aware sound source localization via audio-visual scene understanding. \emph{CVPR}, 2025. arXiv:2506.18557.
\bibitem{vyas2025peav} A.~Vyas et al. Pushing the frontier of audiovisual perception with large-scale multimodal correspondence learning. arXiv:2512.19687, 2025.
\bibitem{wiedenbeck1999icons} S.~Wiedenbeck. The use of icons and labels in an end user application program: An empirical study of learning and retention. \emph{Behaviour \& IT} 18(2):68--82, 1999.
\bibitem{wu2025claira} T.-H. Wu, J.~E. Gonzalez, T.~Darrell, D.~M. Chan. CLAIR-A: Leveraging large language models to judge audio captions. \emph{ASRU}, 2025. arXiv:2409.12962.
\bibitem{wu2025flam} Y.~Wu et al. FLAM: Frame-wise language-audio modeling. \emph{ICML}, 2025. arXiv:2505.05335.
\bibitem{wu2025qwenimage} C.~Wu et al. Qwen-Image technical report. arXiv:2508.02324, 2025.
\bibitem{xu2025qwen3omni} J.~Xu et al. Qwen3-Omni technical report. arXiv:2509.17765, 2025.
\bibitem{yariv2023audiotoken} G.~Yariv, I.~Gat, L.~Wolf, Y.~Adi, I.~Schwartz. AudioToken: Adaptation of text-conditioned diffusion models for audio-to-image generation. \emph{Interspeech}, 2023. arXiv:2305.13050.
\bibitem{yariv2024tempotokens} G.~Yariv, I.~Gat, S.~Benaim, L.~Wolf, I.~Schwartz, Y.~Adi. Diverse and aligned audio-to-video generation via text-to-video model adaptation. \emph{AAAI}, 2024. arXiv:2309.16429.
\bibitem{yoshida2023offscreen} M.~Yoshida, R.~Togo, T.~Ogawa, M.~Haseyama. Off-screen sound separation based on audio-visual pre-training using binaural audio. \emph{Sensors} 23(9):4540, 2023.
\bibitem{zheng2023judge} L.~Zheng et al. Judging LLM-as-a-judge with MT-Bench and Chatbot Arena. \emph{NeurIPS Datasets and Benchmarks}, 2023. arXiv:2306.05685.
\bibitem{zheng2024pride} C.~Zheng, H.~Zhou, F.~Meng, J.~Zhou, M.~Huang. Large language models are not robust multiple choice selectors. \emph{ICLR}, 2024. arXiv:2309.03882.
\bibitem{zhou2022audiocapeval} Z.~Zhou, Z.~Zhang, X.~Xu, Z.~Xie, M.~Wu, K.~Q. Zhu. Can audio captions be evaluated with image caption metrics? \emph{ICASSP}, 2022.

\end{thebibliography}
